\documentclass[10pt,a4paper]{amsart}
\usepackage[margin=19mm]{geometry}
\usepackage{amsmath,amssymb,mathtools}
\usepackage[scr=rsfs,cal=euler]{mathalfa}
\usepackage{xfrac}
\usepackage[unicode,hidelinks,bookmarks=false]{hyperref}
\newcommand{\ie}{i.e.\ }
\newcommand{\cf}{cf.\ }
\newcommand{\resp}{respectively\ }
\newcommand{\Subsection}{\subsection}
\providecommand{\nobreakdash}{\nobreak\hbox{-}}
\setlength{\emergencystretch}{2em}
\multlinegap0pt
\usepackage{xcolor}
\usepackage[shortlabels]{enumitem}
\usepackage{bm}
\usepackage{amssymb}
\usepackage{tcolorbox}
\newtheorem{lemma}{Lemma}
\newcommand{\lec}[2]{\stackrel{\text{#1}}{#2}}
\newcommand{\norm}[1]{\left\lVert#1\right\rVert}
\newcommand{\sob}[2]{\mathbb{W}^{#1,#2}\left(\Omega; \mathbb{R}\right)}
\newcommand{\sobeq}[2]{\widetilde{\mathbb{W}}^{#1,#2}\left(\Omega; \mathbb{R}\right)}
\title{Coded Computing for Resilient Distributed Computing: A Learning-Theoretic Framework}
\author{Parsa Moradi, Behrooz Tahmasebi, Mohammad Ali Maddah-Ali}
\date{Source: arXiv:2406.00300v3 (CC BY 4.0)}
\begin{document}
\maketitle

\noindent\emph{Source and licence.} Coded Computing for Resilient Distributed Computing: A Learning-Theoretic Framework. Version arXiv:2406.00300v3 (CC BY 4.0), distributed by arXiv under the Creative Commons Attribution 4.0 International licence, http://creativecommons.org/licenses/by/4.0/, as recorded in the arXiv licence metadata for this exact version and shown on the abstract page https://arxiv.org/abs/2406.00300v3. Full licence notice: \texttt{licenses/arxiv-2406.00300-CC-BY-4.0.txt}.\par\smallskip

\noindent\textit{Editorial scope.} Counted unit: the article's lemma lem:norm-equiv (source0.tex lines 1161--1168). The article's complete original proof of it is delivered with it (source0.tex lines 1169--1228, one \texttt{proof} environment of 60 lines). No mathematics is written, repaired or re-derived: the statement and the proof are the article's own lines, in the article's own notation, and no retained line is substituted or re-flowed.  No further source-local context is needed: every cross-reference of every retained line resolves inside the two delivered spans.  Nothing was omitted from any retained span, so the package carries no omission record.  No retained line contains a citation, so the wrapper carries no bibliography and no bibliography entry.  No retained line contains a figure environment, an image-inclusion command or drawing code, so no image is delivered and the package declares no asset dependency.  Editorial formatting only, in addition: an AMS \texttt{amsart} preamble that restates the article's own declarations and macros used by the retained lines, namely the environments \texttt{lemma} on separate counters, together with the standard packages those lines need.  The retained environments print in this document as Lemma 1 in order of appearance, whereas the article prints the same items with the numbering of its own full document.  The complete native source of the article as distributed in the arXiv e-print for this exact version is bundled unchanged as \texttt{sources/160100/source0.tex}.
\begin{lemma}\label{lem:norm-equiv}
    Let $\Omega = (a,b)$ be an arbitrary open interval in $\mathbb{R}$. Then for every $g \in \sob{2}{2}$:
    \begin{align}
        \label{lem:norm-equiv1}
        &\norm{g}^2_{\sob{2}{2}}\leqslant \left[2(b-a)\max\{1, (b-a)\} \left (2\max\{1, (b-a)\}^2 + 1 \right) +1\right] \cdot\norm{g}^2_{\sobeq{2}{2}} \\ \label{lem:norm-equiv2}
        &\norm{g}^2_{\sobeq{2}{2}}\leqslant \left(\frac{4}{(b-a)}\max\{1, (b-a)\}^2 + 1\right)\cdot \norm{g}^2_{\sob{2}{2}}.
    \end{align}
\end{lemma}

\begin{proof}
    By expanding $g$ around the point $a$ and utilizing the integral remainder form of Taylor's expansion, for every $x \in (a,b)$ we have:
    \begin{align}
        g(x) = g(a) + \int_a^x g'(t)\,dt.
    \end{align}
    Therefore: \begin{align}\label{eq:semi_norm_neq1}
        g(x)^2 &= g(a)^2 + \left(\int_a^x g'(t)\,dt \right)^2 + 2g(a)\cdot\int_a^x g'(t)\,dt \nonumber \\ 
        &\lec{(a)}{\leqslant} 2g(a)^2 + 2\left(\int_a^x g'(t)\,dt \right)^2 \nonumber \\ 
        &\lec{(b)}{\leqslant} 2g(a)^2 + 2(x-a)\int_a^x g'(t)^2\,dt \nonumber \\  
        &\lec{(c)}{\leqslant}  2g(a)^2 + 2(b-a)\int_a^b g'(t)^2\,dt,
    \end{align}
    where (a) follows from the $ 2g(a)\cdot\int_a^x g'(t)\,dt \leqslant g(a)^2 + \left(\int_a^x g'(t)\,dt \right)^2$, (b) is based on Cauchy-Schwartz inequality, and (c) is due to $x \leqslant b$. Following the same steps as \eqref{eq:semi_norm_neq1}, we have:
\begin{align}\label{eq:semi_norm_neq2}
        g'(x)^2 \leqslant 2g'(a)^2 + 2(b-a)\int_a^b g''(t)^2\,dt.
    \end{align}
    Integrating both sides of \eqref{eq:semi_norm_neq2} we have:
\begin{align}\label{eq:semi_norm_neq3}
        \int_a^b g'(y)^2\,dy &\leqslant 2\int_a^b g'(a)^2\,dy + 2(b-a)\int_a^b \left(\int_a^b g''(t)^2\,dt\right) \,dy \nonumber \\ 
        &= 2(b-a).g'(a)^2 + 2(b-a)^2 \int_a^b g''(t)^2\,dt \nonumber \\
        &\lec{(a)}{\leqslant} 2(b-a)\cdot \max\{1, (b-a)\} \cdot \left(g(a)^2 + g'(a)^2 + \int_a^b g''(t)^2\,dt \right)\nonumber \\
        &= 2(b-a)\cdot\max\{1, (b-a)\}\cdot\norm{g}^2_{\sobeq{2}{2}},
    \end{align}
    where (a) is based on adding the positive term $(b-a)\cdot g(a)^2$ to the right-hand side. Based on the \eqref{eq:semi_norm_neq1} and \eqref{eq:semi_norm_neq3} we have:
\begin{align}\label{eq:semi_norm_neq4}
        \int_a^b g(x)^2\,dx &\lec{(a)}{\leqslant} 2\int_a^b g(a)^2\,dx + 2(b-a)\int_a^b \left(\int_a^b g'(t)^2\,dt\right)\,dx \nonumber \\ 
        &= 2(b-a).g(a)^2 + 2(b-a)^2 \int_a^b g'(t)^2\,dt \nonumber \\ 
        &\lec{(b)}{\leqslant} 2(b-a)\cdot g(a)^2 + 4(b-a)^3\cdot g'(a)^2 + 4(b-a)^4 \int_a^b g''(t)^2\,dt \nonumber \\ 
        &\leqslant 4(b-a)\cdot\max\{1, (b-a)^3\}\cdot\norm{g}^2_{\sobeq{2}{2}},
    \end{align}
    where (a) follows by integrating both sides of \eqref{eq:semi_norm_neq1} and (b) is due to \eqref{eq:semi_norm_neq2}. Using \eqref{eq:semi_norm_neq4} and \eqref{eq:semi_norm_neq3} and the fact that $\int_a^b g''(t)^2\, dt \leqslant \|g\|^2_{\sob{2}{2}}$, we have:
    \begin{align}
        \norm{g}^2_{\sob{2}{2}} &= \int_a^b g''(t)^2\, dt +  \int_a^b g'(t)^2\, dt + \int_a^b g(t)^2\, dt \nonumber \\  &\leqslant \left[2(b-a)\max\{1, (b-a)\} \left (2\max\{1, (b-a)\}^2 + 1 \right) +1\right] \cdot\norm{g}^2_{\sobeq{2}{2}}.
    \end{align}
    For the other side, using the same steps as in \eqref{eq:semi_norm_neq1}, we have:
    \begin{align}\label{eq:semi_norm_neq5}
        g(a)^2 &= g(x)^2 + \biggl(\int_a^x g'(t)\,dt \biggr)^2 - 2g(a).\int_a^x g'(t)\,dt \nonumber \\ 
        &\lec{(a)}{\leqslant} 2g(x)^2 + 2\biggl(\int_a^x g'(t)\,dt \biggr)^2 \nonumber \\
        &\lec{(b)}{\leqslant} 2g(x)^2 + 2(x-a)\int_a^x g'(t)^2\,dt \nonumber \\  
        &\lec{(c)}{\leqslant}  2g(x)^2 + 2(b-a)\int_a^b g'(t)^2\,dt,
    \end{align}
    where (a) is because of $-2g(a)\cdot\int_a^x g'(t)\,dt \leqslant g(a)^2 + \left(\int_a^x g'(t)\,dt \right)^2$, (b) follows from Cauchy-Schwartz inequality, and (c) is due to $x \leqslant b$. Integrating both sides of \eqref{eq:semi_norm_neq5}, we have
    \begin{align}\label{eq:semi_norm_neq6}
         (b-a)\cdot g(a)^2 = \int_a^b g(a)^2\,dx  &\leqslant 2\int_a^b g(x)^2\,dx + 2(b-a)\int_a^b \left(\int_a^b g'(t)^2\,dt \right)\,dx \nonumber \\ 
        &\lec{(a)}{=} 2\int_a^b g(t)^2\,dt + 2(b-a)^2 \int_a^b g'(t)^2\,dt \nonumber \\
        &\lec{(b)}{\leqslant} 2\max\{1, (b-a)\}^2\left(\int_a^b g(t)^2\,dt + \int_a^b g'(t)^2\,dt + \int_a^b g''(t)^2\,dt\right) \nonumber \\
        &\leqslant 2\max\{1, (b-a)\}^2 \norm{g}^2_{\sob{2}{2}},
    \end{align}
    where (a) follows by a change of variable from $x$ to $t$ and (b) follows by adding the positive term $2\max\{1, (b-a)\}^2 \cdot \int_a^b g''(t)^2\,dt$ to the right side. Thus, \eqref{eq:semi_norm_neq6} directly results in
    \begin{align}\label{eq:semi_norm_neq7}
        g(a)^2 \leqslant \frac{2}{(b-a)}\max\{1, (b-a)\}^2 \norm{g}^2_{\sob{2}{2}}.
    \end{align}
    Following same steps as \eqref{eq:semi_norm_neq6} and \eqref{eq:semi_norm_neq7} results in
    \begin{align}
        g'(a)^2 \leqslant \frac{2}{(b-a)}\max\{1, (b-a)\}^2 \norm{g}^2_{\sob{2}{2}}.
    \end{align}
    Considering the fact that $\int_a^b g''(t)^2\,dt \leqslant \norm{g}^2_{\sob{2}{2}}$, we can complete the proof:
    \begin{align}
        \norm{g}^2_{\sobeq{2}{2}} &= g(a)^2 + g'(a)^2 + \int_a^b g''(t)^2\,dt  \nonumber \\ &\leqslant \left(\frac{4}{(b-a)}\max\{1, (b-a)\}^2 + 1\right)\cdot\norm{g}^2_{\sob{2}{2}}.
    \end{align}
\end{proof}

\end{document}
